\documentclass[11pt]{article}
\usepackage[T1]{fontenc}
\usepackage{lmodern}
\usepackage{microtype}
\usepackage[margin=1in]{geometry}
\usepackage{amsmath,amssymb,amsthm,mathtools}
\usepackage{booktabs,longtable,tabularx,array}
\usepackage{enumitem}
\usepackage{hyperref}
\usepackage{xurl}
\hypersetup{colorlinks=true,linkcolor=blue,citecolor=blue,urlcolor=blue,pdftitle={A Source-Bound Arithmetic Replay at Radius 0.113231 for the 100-Disk Covering Problem},pdfauthor={}}
\setlist{nosep}
\setlength{\emergencystretch}{2em}
\newtheorem{theorem}{Theorem}
\newtheorem{proposition}{Proposition}
\newtheorem{lemma}{Lemma}
\newtheorem{definition}{Definition}
\newcommand{\D}{\mathbb D}
\newcommand{\R}{\mathbb R}
\newcommand{\rd}{r_D}
\newcommand{\cR}{R_0}
\newcommand{\Q}{\mathbb Q}
\title{A Source-Bound Arithmetic Replay at Radius $0.113231$ for the 100-Disk Covering Problem}
\date{}
\begin{document}
\maketitle
\begin{abstract}
Let $\rd(100)$ denote the least radius of 100 congruent closed disks whose union covers the closed unit disk. We document a source-bound replay, at the exact rational radius $\cR=113231/10^6$, of a six-branch certificate for the lower-bound argument. A fresh transformed copy of a fixed release was checked file-by-file; only occurrences of the exact baseline radius were changed. The complete arithmetic entry was run twice, and a separate fresh-process replay recorded verifier parameters and source hashes. All six boundary-count branches and the 96-case reduction excluding at most 99 effective sites returned positive exact rational margins. The largest-boundary branch visited 3,767,981 rational boxes. The replay consumed no saved branch PASS reports at its acceptance boundary and no certificate at a larger radius. This is a computational verification of the finite arithmetic layer of an existing geometric argument, not a proof-assistant formalization or an independent rederivation of its geometric lemmas. We state the lower-bound conclusion only conditionally on the normalization, planar-diagram, support, cancellation, degeneracy, and compactness steps specified below.
\end{abstract}

\noindent\textbf{Keywords:} disk covering; exact rational arithmetic; reproducible computation; source integrity; geometric optimization.

\section{Problem and scope}
Let $\D=\{x\in\R^2:\|x\|\leq 1\}$. Define
\[
 \rd(n)=\inf\{r\geq0:\text{there exist }c_1,\ldots,c_n\in\R^2\text{ with }\D\subseteq\bigcup_{i=1}^n\overline B(c_i,r)\}.
\]
The centers are unrestricted and disks are closed. This is the classical circular covering problem, studied computationally by Zahn and in later work on covering constants; see \cite{Zahn1962,Kershner1939,FejesToth2005}.

The purpose of this paper is narrower than a new geometric solution. We give a reproducible, exact-radius computational replay at
\[
 \cR=\frac{113231}{1000000}=0.113231.
\]
The source release already contains a geometric proof architecture and executable arithmetic verifiers. The work reported here specializes its exact radius constants in an isolated copy, executes the complete lower-bound entry from the bottom up, and independently checks source and runtime dependencies. It does not supply a new formal proof of every geometric implication used by the source release.

\begin{theorem}[Conditional certificate statement]\label{thm:conditional}
Assume the geometric transfer package in Section~\ref{sec:geometry}: (G1) radius-$\cR$ center normalization preserves coverage and permits construction of one common clipped diagram; (G2) the diagram's boundary/interior count identities and the reduction to the 96 cases with at most 99 effective sites are valid; (G3) the listed star, face, hinge, signed-edge, and capped-pair inequalities bound the relevant area contributions for every admissible diagram; (G4) the planar, core-reduction, and degeneracy arguments cover all limiting cases; and (G5) the enlargement/compactness/continuity argument transfers strict exclusion at the endpoint to a lower bound for the infimum. Then the replayed finite arithmetic certificate implies
\[
 \rd(100)>\frac{113231}{1000000}.
\]
The computational result established here is that the exact rational predicates and finite enumerations used by this certificate pass at $\cR$, under the source hashes and run conditions recorded in Section~\ref{sec:replay}. The geometric assumptions are not formalized by the replay software.
\end{theorem}

The distinction in Theorem~\ref{thm:conditional} is essential: the output `PASS' is a statement about executable arithmetic checks, not by itself about the truth of the continuum geometric transfer package. In particular, neither an exact optimum nor a lower bound at $0.11324$ is asserted here.

\section{Normalization and finite classification}
\label{sec:architecture}
A crucial feature of the source proof is that it normalizes all centers at the same target radius before rebuilding the diagram. Put
\[
 C=\sqrt{1-\cR^2}=\sqrt{\frac{987178740639}{10^{12}}}.
\]
Suppose a center is $du$ with $d>C$ and $\|u\|=1$. Replace it by $d'u$, where $d'=(1-\cR^2)/d$. For every $z\in\D$,
\begin{equation}\label{eq:normalize}
 \|z-d'u\|^2-\cR^2=\frac{d'}d\bigl(\|z-du\|^2-\cR^2\bigr)+\left(1-\frac{d'}d\right)(\|z\|^2-1).
\end{equation}
Here $0<d'/d<1$. Thus if $z$ was covered by the old disk, the first term on the right is nonpositive, and the second term is nonpositive because $z\in\D$. Hence $z$ remains covered. Applying this operation to every center yields a representative with $\|c_i\|\leq C$. The diagram and all counts are then formed anew; no old ownership or boundary-count information is assumed to persist.

The exact squared norm cap used in the run is
\[
 C^2=1-\cR^2=\frac{987178740639}{10^{12}}.
\]
The global arithmetic entry first excludes configurations with at most 99 effective sites. It tests 96 count cases; the tightest reported case has $B=34$ and an exact rational minimum margin. Its decimal display is $5.48874303\\times 10^{-4}>0$; the audit checks the sign from the exact certificate value.
\[
 5.48874303\times 10^{-4}>0.
\]
After this reduction, the surviving boundary counts in the common $m=100$ diagram are $B=31,\ldots,39$. They are partitioned into six arithmetic branches: $B=31,32,33,34,35$, and $B\geq36$.

\section{The six arithmetic branches}
\label{sec:branches}
Each component verifier returns a rational lower bound for its minimum final area margin. The table reports decimal displays rounded for readability; acceptance in the code uses the exact rational values stored in the cited JSON certificate. The exact rational for every row is checked by the companion audit script.
\begin{center}
\begin{tabular}{@{}lll@{}}
\toprule
Branch & Scope of the finite check & Positive margin (display)\\
\midrule
$B=31$ & Global curvature rows; full-face and vertex supports & $0.003738017486$\\
$B=32$ & Low-$M$ profiles and high-$M$ core split & $0.000267991536$\\
$B=33$ & Coherent hinge and metric rows & $0.000090870200$\\
$B=34$ & Single hinge-bank rows & $0.000179345825$\\
$B=35$ & Low-degree/ear reductions and exceptional $M=26$ core & $0.000025324509$\\
$B\geq36$ & Capped-pair tail over the full declared domain & $0.000290874695$\\
\bottomrule
\end{tabular}
\end{center}
All six entries were run twice at the same exact radius. The second run compared each branch's exact margin with the first run. For $B\geq36$, the rational branch-and-bound visited 3,767,981 boxes and accepted 1,879,204; its exact parameters include $\lambda=3/50$, $\mu=92531/10000000$, and target $751/100000$. These counts describe the recorded computation and do not replace the correctness of the box enclosure rule.

\subsection{What the branches check}
The branch labels denote the boundary count in the rebuilt diagram, not a partition of an inherited graph. In outline, the $B=31$ branch combines global degree/curvature restrictions with full-face inequalities; $B=32$ combines a direct low-$M$ profile check with high-$M$ core reductions; $B=33$ uses coherent replacement rows and high-$M$ metric cases; $B=34$ uses a single boundary hinge bank; $B=35$ combines low-degree and exceptional-core reductions; and $B\geq36$ is handled by the capped-pair tail. Detailed source arguments and case definitions are retained in the fixed release package. The executable certificate verifies their finite rational consequences, including the source assertions, but is not a proof assistant checking the geometric derivations.

\section{Replay protocol and provenance}
\label{sec:replay}
The source package is the fixed `certified\_bounds\_11323' release associated with upstream commit `67ebf1c5776f72a7b7786a8655eed5463429ae55'~\\cite{Disk100Release}. A fresh isolated copy was made. The radius conversion changed only exact occurrences of $11323/100000$ and `F(11323,100000)' into $113231/1000000$ and its corresponding rational string. The run records 626 source files, 60 transformed files, 158 Python files parsed, and 1,662 assertions preserved. The source and transformed-copy snapshots matched their manifests after execution.

The full entry was run twice: the initial fresh run took about 260.12 seconds; the separately instrumented replay took about 781.49 seconds. The latter recorded 20 processes and 28 loaded package sources, with no verifier entry carrying a different global radius. Each observed branch radius equals $113231/1000000$. A full-file manifest binds dynamically loaded sources that are not visible as ordinary imports to the recorder. Negative controls rejected an optimized/assertions-disabled run and a verifier deliberately passed the wrong radius; both exited with code 1.

No saved branch PASS report was used as acceptance evidence, and the $0.113235$ certificates were not used. The direct entry reran its six branches and repeated radius/margin assertions. These are source-bound replay and dependency claims, rather than claims that every Python frame was instrumented.

\section{Geometric transfer obligations}
\label{sec:geometry}
The finite certificate is useful only together with the following mathematical obligations. They are stated explicitly because they are not discharged by the process trace or the rational arithmetic audit.
\begin{enumerate}[label=(G\arabic*)]
 \item \textbf{Same-radius normalization and diagram.} The transformation in \eqref{eq:normalize} preserves coverage for the fixed endpoint radius. The geometric proof must justify construction of a common clipped diagram after normalizing all centers and must not transfer ownership or counts from a prior diagram.
 \item \textbf{Count reduction.} The boundary/interior identities, effective-site reductions, and exclusion of all 96 cases with at most 99 effective sites must apply to arbitrary admissible covers, not only generic diagrams.
 \item \textbf{Local area inequalities.} Star, face, hinge, and metric supports must bound actual geometric area for every allowed label/profile. The charged signed-edge terms must cancel with the stated multiplicities, without dropping positive interface contributions or counting a physical edge more than once.
 \item \textbf{Planarity and degeneracies.} The planar graph and combinatorial isoperimetric inequalities, core reductions, limiting configurations, ties, boundary coincidences, and zero-area faces require proofs covering degenerate cases.
 \item \textbf{Strict endpoint.} The generic perturbation/enlargement procedure must preserve the required bounds uniformly as the radius decreases to $\cR$. Compactness of the projected center space and continuity of the covering-radius functional then support attainment and the strict endpoint conclusion.
\end{enumerate}
The source release gives analytic arguments for these steps, but this paper's independent work verifies the transformed numerical chain and its provenance rather than re-proving those arguments. A reader accepting (G1)--(G5) obtains the stated strict lower bound from the finite certificate; without that acceptance, the unconditional geometric conclusion remains open here.

\section{Limitations and interpretation}
The work establishes a reproducible arithmetic specialization of an existing proof chain at a new exact radius. Its evidential strengths are exact rational assertions, a complete six-branch entry executed twice, explicit failure controls, and source/module hash binding. Its limitations are equally concrete:
\begin{itemize}
 \item no proof assistant formalization of the analytic geometric lemmas;
 \item no independent implementation of the entire geometry-to-area argument;
 \item no exact value or optimal configuration for $\rd(100)$;
 \item no claim at any radius larger than $113231/1000000$.
\end{itemize}
The reported minimum margins are certificates for the specified finite inequalities under their stated premises. They are not direct geometric measurements of arbitrary covers.

\section{Reproducibility}
The source bundle, transformation manifest, raw replay outputs, process traces, controls, and exact certificates are part of the accompanying repository at \url{https://github.com/cjt42554-crypto/disk100-covering-certificates}. The primary machine-readable summary is `verify/out/n100-direct-113231-complete.json'; the independent dependency audit is `verify/out/n100-radius-dependency-audit-direct-113231-fresh.json'; the source transformation record is `runtime/lower-small-advance-direct-113231-fresh/transformation.json'. The paper bundle contains `audit/verify\_claims.py', which validates the numerical claims and hashes against those committed artifacts.

\bibliographystyle{plain}
\bibliography{references}
\appendix
\section{Evidence ledger}
\label{app:ledger}
Table~\ref{tab:ledger} states what the software verifies and what remains a mathematical dependency.
\begin{longtable}{@{}p{0.25\linewidth}p{0.27\linewidth}p{0.39\linewidth}@{}}
\caption{Claim-to-evidence ledger.}\label{tab:ledger}\\
\toprule Claim & Evidence status & Independent check / limitation\\
\midrule
\endfirsthead
\toprule Claim & Evidence status & Independent check / limitation\\
\midrule
\endhead
Exact target radius $113231/10^6$ & Confirmed & Rational parsing and equality in both complete outputs and all branch rows.\\
Six branches have positive margins & Confirmed arithmetically & Two full runs; exact Fraction comparison and repeat equality; geometric meanings depend on (G3).\\
At most 99 effective sites excluded & Confirmed arithmetically & 96 finite cases and positive exact minimum margin; geometric case coverage depends on (G1)--(G2).\\
No higher-radius PASS inherited & Confirmed by entry and certificate metadata & Full fresh entry; explicit flags state no saved PASS acceptance and no 0.113235 evidence.\\
Source transformation constrained & Confirmed by hash audit & 626 files snapshotted, 60 changed by exact radius substitutions; 1,662 assertions preserved.\\
Runtime radius/source binding & Confirmed within recorder scope & 20 processes, 28 loaded package sources; full manifest additionally binds dynamic files. Not every Python frame is instrumented.\\
Geometric transfer from inequalities to all covers & Not independently verified here & Requires (G1)--(G5), including degeneracies and endpoint compactness.\\
Exact optimum $\rd(100)$ & Unverified & No matching upper/lower endpoint or optimality proof is claimed.\\
\bottomrule
\end{longtable}

\section{Exact branch margins}
For convenient machine checking, the exact branch margins are retained verbatim in `verify/out/n100-direct-113231-complete.json'. In branch order $31,32,33,34,35,\geq36$, they are
\begin{align*}
&\frac{1869008743124689739451}{500000000000000000000000},\quad
\frac{53598307165582960111}{200000000000000000000000},\\
&\frac{90870199742281553043}{1000000000000000000000000},\quad
\frac{179345824655025580389}{1000000000000000000000000},\\
&\frac{633112725883465801}{25000000000000000000000},\quad
\frac{29087469503098483006469778967931605940457120574226219}{100000000000000000000000000000000000000000000000000000000}.
\end{align*}
The decimals in Section~\ref{sec:branches} are rounded representations only; these rational values are the acceptance values.
\end{document}
